\documentclass[exercice]{thedocument}%
\startdatakeys
\source{}
\cycle{}
\competencesDuSocle{}
\connaissancesRequises{}
\competencesTravaillees{}
\enddatakeys
\begin{document}%
Pour fêter les 25 ans de sa boutique, un chocolatier souhaite offrir aux premiers clients de la journée une boîte contenant des truffes au chocolat.
\medskip
Il a confectionné $300$ truffes: $125$ truffes parfumées au café et $175$ truffes enrobées de noix de coco. Il souhaite fabriquer ces boîtes de sorte que:
\medskip
\begingroup\parindent=1cm
$\bullet$ Le nombre de truffes parfumées au café soit le même dans chaque boîte;
$\bullet$ Le nombre de truffes enrobées de noix de coco soit le même dans chaque boîte;
$\bullet$ Toutes les truffes soient utilisées.
\endgroup
\begin{enumerate}
    \item Décomposer $125$ et $175$ en produit de facteurs premiers.
    \item En déduire la liste des diviseurs communs à $125$ et $175$.
    \item Quel nombre maximal de boîtes pourra-t-il réaliser ?
    \item Dans ce cas, combien y aura-t-il de truffes de chaque sorte dans chaque boîte?
\end{enumerate}
\end{document}
